\documentclass[12pt,a4paper]{ctexart}

% ==================== 宏包引用 ====================
\usepackage[left=2.0cm, right=2.0cm, top=2.5cm, bottom=2.5cm]{geometry} % 页边距
\usepackage{amsmath, amssymb} % 数学公式
\usepackage{siunitx} % 单位符号
\usepackage{tikz} % 绘图
\usetikzlibrary{calc, arrows.meta, shapes.geometric}
\usepackage[table]{xcolor} % 颜色 (启用table选项支持\rowcolor)
\usepackage{enumitem} % 列表定制
\usepackage{caption} % 图注
\usepackage{array} % 表格增强
\usepackage{fancyhdr} % 页眉页脚

% ==================== 样式设置 ====================

% 定义书本颜色
\definecolor{bookblue}{RGB}{0, 155, 230}    % 标题蓝
\definecolor{weekendred}{RGB}{200, 50, 50}  % 周末红
\definecolor{solargreen}{RGB}{50, 150, 50}  % 节气绿
\definecolor{grassgreen}{RGB}{196, 214, 160} % 草绿色

% 页面风格设置
\pagestyle{fancy}
\fancyhf{} % 清空页眉页脚
\renewcommand{\headrulewidth}{0pt} % 去掉页眉线

% 去掉默认段首缩进，改为手动控制
\setlength{\parindent}{0pt}

% 自定义命令：农历小字
\newcommand{\lunar}[1]{\scalebox{0.7}{#1}}

\begin{document}

% ==================== 第 53 页内容 ====================

% --- 顶部标题栏 ---
\begin{tikzpicture}[remember picture, overlay]
    % 蓝色顶栏装饰
    \fill[bookblue!10] (current page.north west) rectangle ($(current page.north east)+(0,-1.5)$);
    \node[anchor=west, bookblue] at ($(current page.north west)+(2.5,-1)$) {\Large \textbf{习题 2.5}};
\end{tikzpicture}

\vspace{0.5cm}

\noindent
{\color{bookblue}\large\textbf{A 组}}

\vspace{0.5cm}

% --- 第 1 题 ---
\textbf{1.} 某市政府为落实保障性住房政策，去年已投入 3 亿元资金，并规划投入资金逐年增加，明年将投入 12 亿元资金用于保障性住房建设. 求这两年中投入资金的年平均增长率.

\vspace{0.8cm}

% --- 第 2 题 ---
\textbf{2.} 某店只销售某种进价为 40 元/kg 的特产. 已知该店按 60 元/kg 出售时，平均每天可售出 100 kg，后来经过市场调查发现，单价每降低 1 元，则平均每天的销售量可增加 10 kg. 若该店销售这种特产计划平均每天获利 2\,240 元.
\begin{enumerate}[label=(\arabic*), leftmargin=3.5em, itemsep=0.2em]
    \item 每千克该种特产应降价多少元？
    \item 为尽可能让利于顾客，则该店应按原售价的几折出售？
\end{enumerate}

\vspace{0.8cm}

% --- 第 3 题 ---
\noindent
\begin{minipage}[t]{0.52\textwidth}
    \textbf{3.} 某单位准备将院内一块长 30 m、宽 20 m 的长方形空地，建成一个矩形花园，要求在花园中修建两条纵向和一条横向的小道，剩余的地方种植花草，如图所示. 要使种植花草的面积为 532 m$^2$，那么小道进出口的宽度应为多少米？
    
    \smallskip
    (注：所有小道进出口的宽度相等.)
\end{minipage}%
\hfill
\begin{minipage}[t]{0.45\textwidth}
    \centering
    \begin{tikzpicture}[scale=0.13]
        % 定义参数
        \def\W{30}
        \def\H{20}
        \def\pw{2} % 示意小道宽
        
        % 标注线样式
        \tikzstyle{dim} = [<->, >=latex, thin]
        
        % 绘制草地底色 (绿色块)
        \fill[grassgreen] (0,0) rectangle (\W, \H);
        
        % 绘制小道 (灰白填充 + 黑框)
        % 横向
        \filldraw[fill=black!5, draw=black] (0, \H/2 - \pw/2) rectangle (\W, \H/2 + \pw/2);
        % 纵向1 (左)
        \filldraw[fill=black!5, draw=black] (\W/3 - \pw/2, 0) rectangle (\W/3 + \pw/2, \H);
        % 纵向2 (右)
        \filldraw[fill=black!5, draw=black] (2*\W/3 - \pw/2, 0) rectangle (2*\W/3 + \pw/2, \H);
        
        % 外框重描
        \draw[thick] (0,0) rectangle (\W, \H);
        
        % 尺寸标注
        \draw[dim] (0, -2) -- (\W, -2) node[midway, fill=white] {30 m};
        % \draw[dim] (\W+2, 0) -- (\W+2, \H) node[midway, fill=white] {20 m};
        % \draw[thin] (0,0) -- (0, -2.5);
        % \draw[thin] (\W,0) -- (\W, -2.5);
        % \draw[thin] (\W,0) -- (\W+2.5, 0);
        % \draw[thin] (\W,\H) -- (\W+2.5, \H);
        
        % 纵向箭头 (右侧) - 模拟原图的蓝色箭头
        \draw[bookblue, ->, >=stealth] (\W+1, \H) -- (\W+1, \H/2 + 2);
        \draw[bookblue, ->, >=stealth] (\W+1, 0) -- (\W+1, \H/2 - 2);
        \node at (\W+1, \H/2) {\scriptsize \color{bookblue}20 m};
        
        \node at (\W/2, -5) {(第 3 题图)};
    \end{tikzpicture}
\end{minipage}

\vspace{0.5cm}

% --- 第 4 题 ---
\textbf{4.} 如图是某年 8 月的日历表，在此日历表上可以用一个矩形圈出 3$\times$3 个位

\begin{center}
    \begin{minipage}{0.6\textwidth}
        \centering
        \renewcommand{\arraystretch}{1.3}
        \setlength{\tabcolsep}{5pt}
        % 日历表格
        \begin{tabular}{|c|c|c|c|c|c|c|}
            \hline
            \rowcolor{black!5} % 表头背景色
            \textcolor{weekendred}{日} & 一 & 二 & 三 & 四 & 五 & \textcolor{weekendred}{六} \\
            \hline
             & & & 1 & 2 & 3 & \textcolor{weekendred}{4} \\
             & & & \lunar{建军节} & \lunar{十五} & \lunar{十六} & \textcolor{weekendred}{\lunar{十七}} \\
            \hline
            \textcolor{weekendred}{5} & 6 & 7 & 8 & 9 & 10 & \textcolor{weekendred}{11} \\
            \textcolor{weekendred}{\lunar{十八}} & \lunar{十九} & \textcolor{solargreen}{\lunar{立秋}} & \lunar{廿一} & \lunar{廿二} & \lunar{廿三} & \textcolor{weekendred}{\lunar{廿四}} \\
            \hline
            \textcolor{weekendred}{12} & 13 & 14 & 15 & 16 & 17 & \textcolor{weekendred}{18} \\
            \textcolor{weekendred}{\lunar{廿五}} & \lunar{廿六} & \lunar{廿七} & \lunar{廿八} & \lunar{廿九} & \lunar{七月} & \textcolor{weekendred}{\lunar{初二}} \\
            \hline
            \textcolor{weekendred}{19} & 20 & 21 & 22 & 23 & 24 & \textcolor{weekendred}{25} \\
            \textcolor{weekendred}{\lunar{初三}} & \lunar{初四} & \lunar{初五} & \lunar{初六} & \lunar{初七} & \lunar{初八} & \textcolor{weekendred}{\lunar{初九}} \\
            \hline
            \textcolor{weekendred}{26} & 27 & 28 & 29 & 30 & 31 & \\
            \textcolor{weekendred}{\lunar{初十}} & \lunar{十一} & \lunar{十二} & \lunar{十三} & \lunar{十四} & \lunar{十五} & \\
            \hline
        \end{tabular}
        \smallskip
        \small{(第 4 题图)}
    \end{minipage}
\end{center}

% 53页页脚
\vfill
\begin{tikzpicture}[remember picture, overlay]
    \node[anchor=south east] at ($(current page.south east)+(-2,1)$) {
        \small \textcolor{orange!70!black}{第 2 章} \quad \textcolor{gray}{一元二次方程} \quad \textbf{\color{orange!70!black}\large 53}
    };
\end{tikzpicture}

\newpage

% ==================== 第 54 页内容 ====================

% --- 第 7 题 ---
\noindent
\begin{minipage}[t]{0.50\textwidth}
    \textbf{7.} 如图，在矩形 $ABCD$ 中， $BC=24$ cm， $P, Q, M, N$ 分别从点 $A, B, C, D$ 同时出发，分别沿边 $AD, BC, CB, DA$ 移动，且当有一个点先到达所在边的另一个端点时，其他各点也随之停止移动. 已知移动一段时间后，若 $BQ=x$ cm ($x \neq 0$)，则 $AP=2x$ cm， $CM=3x$ cm， $DN=x^2$ cm.
\end{minipage}%
\hfill
\begin{minipage}[t]{0.48\textwidth}
    \centering
    \begin{tikzpicture}[scale=0.28]
        % 坐标定义
        \coordinate (A) at (0, 8);
        \coordinate (B) at (0, 0);
        \coordinate (C) at (16, 0);
        \coordinate (D) at (16, 8);
        
        % 绘制矩形
        \draw[thick] (A) -- (B) -- (C) -- (D) -- cycle;
        
        % 标注顶点
        \node[left] at (A) {$A$};
        \node[left] at (B) {$B$};
        \node[right] at (C) {$C$};
        \node[right] at (D) {$D$};
        
        % 动点位置 (示意)
        \coordinate (P) at (3, 8);
        \coordinate (N) at (13, 8);
        \coordinate (Q) at (2, 0);
        \coordinate (M) at (11, 0);
        
        % 连线
        \draw[thick] (P) -- (Q);
        \draw[thick] (N) -- (M);
        
        % 动点标注
        \fill (P) circle (3pt) node[above] {$P$};
        \fill (N) circle (3pt) node[above] {$N$};
        \fill (Q) circle (3pt) node[below] {$Q$};
        \fill (M) circle (3pt) node[below] {$M$};
        
        % 运动方向箭头 (蓝色)
        \draw[->, >=stealth, bookblue, thick] (A) -- (1.5, 8);
        \draw[->, >=stealth, bookblue, thick] (D) -- (14.5, 8);
        \draw[->, >=stealth, bookblue, thick] (B) -- (1, 0);
        \draw[->, >=stealth, bookblue, thick] (C) -- (13.5, 0);
        
        \node at (8, -2.5) {(第 7 题图)};
    \end{tikzpicture}
\end{minipage}

\begin{enumerate}[label=(\arabic*), leftmargin=3.5em, itemsep=0.4em]
    \item 当 $x$ 为何值时， $P, N$ 两点重合？
    \item 问 $Q, M$ 两点能重合吗？若 $Q, M$ 两点能重合，则求出相应的 $x$ 的值；若 $Q, M$ 两点不能重合，请说明理由.
    \item 当 $x$ 为何值时，以 $P, Q, M, N$ 为顶点的四边形是平行四边形？
\end{enumerate}

% 54页页脚
\vfill
\begin{tikzpicture}[remember picture, overlay]
    \node[anchor=south west] at ($(current page.south west)+(2,1)$) {
        \textbf{\color{orange!70!black}\large 54} \quad \small \textcolor{gray}{数学 \quad 九年级上册}
    };
\end{tikzpicture}

\end{document}